\section{Slices keep logical coordinates}
\label{sec:sequence-sql-slicing}

A slice selects positions \([a,b)\): it includes \(a\) and excludes \(b\).
BioSeqDB uses zero-based positions, so \(S[0]\) is the first symbol. The
length of a valid slice is \(b-a\), and its first position is again zero.
This rule is independent of byte boundaries and of the fragments used to
store \(S\).

The SQL listing gives concrete examples. The type modifier binds the value to
the DNA4 alphabet and 32-bit storage units. Casting the results to text makes
the selected symbols visible.
\begin{lstlisting}[style=seq-nine-sql,
  caption={Zero-based slices and sequence transforms in SQL.},
  label={lst:sequence-sql-examples}]
WITH v(s) AS (
  VALUES ('ACGTAACCGTTA'::bioseq.seq('dna4', 32))
)
SELECT s[2:10]::text AS middle,              -- GTAACCGT
       ((s[2:10])[1:5])::text AS nested,     -- TAAC
       s[-4:]::text AS suffix,              -- GTTA
       bioseq.reverse(s[2:10])::text AS rev, -- TGCCAATG
       bioseq.reverse_complement(s[2:10])::text AS rc
FROM v;                                    -- rc: ACGGTTAC
\end{lstlisting}

The bounds follow a small set of rules.
\begin{itemize}
  \item A missing lower bound means zero. A missing upper bound means the
        sequence length.
  \item A negative bound counts from the end. For length 12, \(-4\) means 8.
        Bounds outside the sequence are clamped to its ends.
  \item An upper bound below the lower bound gives an empty slice.
        Thus \texttt{s[4:2]} has length zero.
  \item An out-of-range element, such as \texttt{s[12]} in this example,
        gives SQL \texttt{NULL}. An explicit null slice bound also gives
        \texttt{NULL}.
\end{itemize}
The implemented subscripting supports one dimension and unit steps. It does
not support assignment through a subscript.

For a forward view \(V=S[a,b)\), view position \(j\) refers to
\(S[a+j]\). A second slice therefore combines its offsets directly:
\[
  V[c,d)=S[a+c,a+d).
\]
The example's \(S[2,10)[1,5)\) is \(S[3,7)\), which reads
\texttt{TAAC}. Neither step needs to round its bounds to storage units.

Reversal changes the mapping, not the chosen symbols' storage. Let
\(W=\operatorname{reverse}(S[a,b))\). Then
\[
  W[j]=S[b-1-j],\qquad
  W[c,d)=\operatorname{reverse}(S[b-d,b-c)).
\]
A reverse-complement view applies the alphabet's complement after this
position mapping. For DNA4, the pairs are \(A\leftrightarrow T\) and
\(C\leftrightarrow G\).

\begin{figure}[H]
\centering
\begin{tikzpicture}[x=1cm,y=1cm,font=\footnotesize]
  \node[anchor=west,font=\bfseries] at (0,4.8)
    {(a) A slice of a slice keeps exact source positions};
  \node[anchor=east] at (2.3,3.9) {$S$};
  \foreach \i/\s in {0/A,1/C,2/G,3/T,4/A,5/A,6/C,7/C,8/G,9/T,10/T,11/A} {
    \node[seq nine cell,fill=black!3] at (2.7+.72*\i,3.9) {\texttt{\s}};
    \node[font=\scriptsize] at (2.7+.72*\i,4.4) {\i};
  }
  \foreach \i/\s in {2/G,3/T,4/A,5/A,6/C,7/C,8/G,9/T}
    \node[seq nine blue] at (2.7+.72*\i,3.9) {\texttt{\s}};
  \foreach \i/\s in {3/T,4/A,5/A,6/C}
    \node[seq nine orange] at (2.7+.72*\i,3.9) {\texttt{\s}};
  \node[anchor=east] at (2.3,2.85) {$V=S[2,10)$};
  \foreach \i/\s in {0/G,1/T,2/A,3/A,4/C,5/C,6/G,7/T} {
    \node[seq nine blue] at (4.14+.72*\i,2.85) {\texttt{\s}};
    \node[font=\scriptsize] at (4.14+.72*\i,2.38) {\i};
  }
  \foreach \i/\s in {1/T,2/A,3/A,4/C}
    \node[seq nine orange] at (4.14+.72*\i,2.85) {\texttt{\s}};
  \draw[seq nine flow,seqNineBlue] (4.14,3.62)--(4.14,3.15);
  \draw[seq nine flow,seqNineBlue] (9.18,3.62)--(9.18,3.15);
  \node[anchor=west,text=seqNineOrange] at (10.35,2.85)
    {$V[1,5)$};

  \node[anchor=west,font=\bfseries] at (0,1.45)
    {(b) Reversal changes which source position comes first};
  \node[anchor=east] at (2.3,.55) {$W$};
  \foreach \i/\s/\p in {0/T/9,1/G/8,2/C/7,3/C/6,4/A/5,5/A/4,6/T/3,7/G/2} {
    \node[seq nine blue] at (4.14+.72*\i,.55) {\texttt{\s}};
    \node[font=\scriptsize] at (4.14+.72*\i,.03) {$j=\i$};
    \node[font=\scriptsize,text=seqNineBlue] at (4.14+.72*\i,1.08) {\p};
  }
  \foreach \i/\s in {1/G,2/C,3/C}
    \node[seq nine orange] at (4.14+.72*\i,.55) {\texttt{\s}};
  \node[anchor=east,font=\scriptsize,text=seqNineBlue] at (3.63,1.08)
    {source position};
  \node[anchor=west,align=left,text=seqNineOrange] at (10.35,.55)
    {$W[1,4)$\\reads $S[6,9)$\\backward: \texttt{GCC}};
\end{tikzpicture}
\caption{Blue cells form the first slice. Orange cells show a second slice.
The source positions above the reversed view decrease while its own positions
increase. The selected intervals are exact even when a boundary falls inside
a packed unit.}
\label{fig:sequence-slice-coordinates}
\end{figure}

This separation matters for long references. A request for the last 100 symbols
can be mapped to the corresponding source interval before any symbols are
read. Stored fragments may be larger than that interval, so the physical read
can still fetch or decompress extra data. Laziness avoids requiring the whole
source to be decoded. It does not make the storage system's read granularity
equal to one symbol.
